\exercisesfor{8}

We suppose that the residual characteristic \(p\) of the field \(K\) is \(>0\) . We denote by \(\mathfrak{o}_{K}\) the ring of the valuation \(v\) of \(K\) .

\begin{enumerate}
\def\labelenumi{\arabic{enumi}.}
\tightlist
\item
  \begin{enumerate}
  \def\labelenumii{(\alph{enumii})}
  \tightlist
  \item
    An element \(\pi\) of \(K\) is called admissible if \(v(\pi) = \theta\) and \(\pi^{p-1} / p \equiv -1 \pmod{\pi o_K}\) .
  \end{enumerate}
\end{enumerate}

(Note that \(\pi^{p-1}/p\in\mathfrak{o}_{K}\) since \(v(\pi)=\theta.\) )

Show that there exist fields K satisfying the conditions of the paragraph and containing an admissible element (adjoin to \(Q_{p}\) a \((p - 1)\) -th root of \(-p\) ).

\begin{enumerate}
\def\labelenumi{(\alph{enumi})}
\setcounter{enumi}{1}
\tightlist
\item
  Let \(\pi\) be an admissible element of K. Consider the following formal power series with coefficients in K:
\end{enumerate}

\[
e _ {\pi} (\mathrm{U}) = \frac {1}{\pi} e (\pi \mathrm{U}) = \sum_ {n = 1} ^ {\infty} \pi^ {n - 1} \mathrm{U} ^ {n} / n!,
\]

\[
l _ {\pi} (\mathrm{U}) = \frac {1}{\pi} l (\pi \mathrm{U}) = \sum_ {n = 1} ^ {\infty} (- \pi) ^ {n - 1} \mathrm{U} ^ {n} / n.
\]

Show that these series are inverses of one another and that their coefficients belong to \(\mathfrak{o}_{K}\) (use Lemma 2).

\begin{enumerate}
\def\labelenumi{(\alph{enumi})}
\setcounter{enumi}{2}
\tightlist
\item
  Let \(\{\mathrm{U},\mathrm{V}\}\) be a set with two elements and let \(\mathrm{H}(\mathrm{U},\mathrm{V})\) be the Hausdorff series. We write
\end{enumerate}

\[
\mathrm{H} _ {\pi} (\mathrm{U}, \mathrm{V}) = \frac {1}{\pi} \mathrm{H} (\pi \mathrm{U}, \pi \mathrm{V}).
\]

Then \(H_{\pi}(U, V) \in \hat{L}_{K}(\{U, V\})\) .

Show that

\[
\mathrm{H} _ {\pi} (\mathrm{U}, \mathrm{V}) = l _ {\pi} (e _ {\pi} (\mathrm{U}) + e _ {\pi} (\mathrm{V}) + \pi e _ {\pi} (\mathrm{U}). e _ {\pi} (\mathrm{V})).
\]

Deduce that \(H_{\pi} \in \hat{L}_{\mathfrak{o}_{K}}(\{\mathrm{U}, \mathrm{V}\})\) , i.e.~that the coefficients of \(H_{\pi}\) belong to \(o_{K}\) ; hence another proof of Proposition 1.

\begin{enumerate}
\def\labelenumi{(\alph{enumi})}
\setcounter{enumi}{3}
\tightlist
\item
  Let \(\tilde{e}_{\pi}, \tilde{l}_{\pi}\) and \(\tilde{\mathbf{H}}_{\pi}\) denote the series obtained by reducing \(e_{\pi}, l_{\pi}\) and \(\mathbf{H}_{\pi}\) modulo \(\pi \mathfrak{o}_{K}\) ; their coefficients belong to the ring \(\mathfrak{o}_{K} / \pi \mathfrak{o}_{K}\) .
\end{enumerate}

Show that \(\tilde{l}_{\pi}(\mathrm{U}) = \mathrm{U} - \mathrm{U}^{p}\) (note that \(v_{p}(n) < (n - 1)\theta\) if \(n \neq 1, p\) ). Deduce that

\[
\tilde {e} _ {\pi} (\mathrm{U}) = \mathrm{U} + \mathrm{U} ^ {p} + \dots + \mathrm{U} ^ {p ^ {n}} + \dots
\]

and that

\[
\tilde {\mathrm{H}} _ {\pi} (\mathrm{U}, \mathrm{V}) = \tilde {e} _ {\pi} (\mathrm{U}) + \tilde {e} _ {\pi} (\mathrm{V}) - (\tilde {e} _ {\pi} (\mathrm{U}) + \tilde {e} _ {\pi} (\mathrm{V})) ^ {p},
\]

or also

\[
\tilde {\mathrm{H}} _ {\pi} (\mathrm{U}, \mathrm{V}) = \mathrm{U} + \mathrm{V} - \Lambda_ {p} \left(\sum_ {n = 0} ^ {\infty} \mathrm{U} ^ {p ^ {n}}, \sum_ {n = 0} ^ {\infty} \mathrm{V} ^ {p ^ {n}}\right),
\]

where \(\Lambda_{p}\) is defined by the formula

\[
\Lambda_ {p} (\mathrm{U}, \mathrm{V}) = (\mathrm{U} + \mathrm{V}) ^ {p} - \mathrm{U} ^ {p} - \mathrm{V} ^ {p}, \quad \text { cf.   Chapter   I,   } \S 1, \text {   Exercise   19. }
\]

In particular, \(\tilde{\mathrm{H}}_{\pi}(\mathrm{U},\mathrm{V})\) belongs to \(\hat{\mathrm{L}}_{\mathbf{F}_{p}}(\{\mathrm{U},\mathrm{V}\})\) and its homogeneous component of degree \(p\) is \(-\Lambda_p(\mathrm{U},\mathrm{V})\) . Then

\[
\begin{array}{c} \tilde {\mathrm{H}} _ {\pi} (\mathrm{U}, \mathrm{V}) = \tilde {\mathrm{H}} _ {\pi} (\mathrm{V}, \mathrm{U}) \\ \tilde {\mathrm{H}} _ {\pi} (\mathrm{U}, \tilde {\mathrm{H}} _ {\pi} (\mathrm{V}, \mathrm{W})) = \tilde {\mathrm{H}} _ {\pi} (\tilde {\mathrm{H}} _ {\pi} (\mathrm{U}, \mathrm{V}), \mathrm{W}) \quad \text { in } \hat {\mathrm{L}} _ {\mathbf {F} _ {p}} (\{\mathrm{U}, \mathrm{V}, \mathrm{W} \}), \end{array}
\]

where W is a third indeterminate.

\begin{enumerate}
\def\labelenumi{(\alph{enumi})}
\setcounter{enumi}{4}
\tightlist
\item
  Let C(U, V) = H(-U, H(-V, H(U, V))) (``Hausdorff commutator''). Then exp(C(U, V)) = exp(-U) exp(-V) exp(U) exp(V). We write
\end{enumerate}

\[
\mathrm{C} _ {\pi} (\mathrm{U}, \mathrm{V}) = \frac {1}{\pi} \mathrm{C} (\pi \mathrm{U}, \pi \mathrm{V}).
\]

Show that the coefficients of \(C_{\pi}(U,V)\) belong to \(\pi o_{K}\) (use the fact that \(\tilde{H}_{\pi}(U,V)=\tilde{H}_{\pi}(V,U)\) ).†

\begin{enumerate}
\def\labelenumi{\arabic{enumi}.}
\setcounter{enumi}{1}
\tightlist
\item
  We know (§ 6, Exercise 3) that if \(n \geqslant 2\) the component of H(U, V) of bidegree \((n, 1)\) is \(\frac{1}{n!} b_n(\mathrm{ad}\mathrm{U})^n (\mathrm{V})\) , where \(b_n\) is the \(n\) -th Bernoulli number.
\end{enumerate}

Deduce from this and Proposition 1 the inequality

\[
v _ {p} (b _ {n} / n!) \leqslant n / (p - 1).
\]

Recover this result by means of the Clausen-von Staudt Theorem (Functions of a Real Variable, Chapter VI, §2, Exercise 6) and show that there is equality if and only if \(S(n) = p - 1\) .

\begin{enumerate}
\def\labelenumi{\arabic{enumi}.}
\setcounter{enumi}{2}
\tightlist
\item
  Suppose that K contains a primitive \(p\) -th root of unity \(w\) . We write \(\pi = w - 1\) . Using the formula
\end{enumerate}

\[
w ^ {i} - 1 = \pi (1 + w + \dots + w ^ {i - 1})
\]

show that \(v(w^{i} - 1) = v(\pi)\) for \(1 \leqslant i \leqslant p - 1\) and that

\[
\frac {w ^ {i} - 1}{\pi} \equiv i \pmod {\pi}.
\]

Deduce, by means of the formula \(p = \prod_{i=1}^{p-1} (w^i - 1)\) , that \(\pi\) is admissible (Exercise 1).

¶ 4. Let c be an integer \(\geqslant1\) and let P be a set of prime numbers containing all prime numbers \(\leqslant c\) . Let \(Z_{P}=S_{P}^{-1}Z\) (§ 4, Exercise 16). Show that the terms of degree \(\leqslant c\) in the Hausdorff series H(U, V) belong to \(L_{Z_{P}}(\{U, V\})\) . Deduce that, if g is a nilpotent Lie \(Z_{P}\) -algebra of class \(\leqslant c\) , the law of composition \((u,v)\mapsto\mathrm{H}(u,v)\) makes g into a P-torsion-free P-divisible nilpotent group of class \(\leqslant c\) . Show that conversely every group with these properties can be obtained by this process. (Use ``the inversion of the Hausdorff formula'', cf.~§6, Exercise 4, and show that the exponents \(\alpha(m)\) , \(\beta(m)\) which appear there belong to \(Z_{P}\) when \(l(m)\leqslant c\) .)

In particular, every p-group of order \(p^{n}\) and class \textless p can be obtained by means of the Hausdorff law from a nilpotent Lie Z-algebra of class \textless p with \(p^{n}\) elements. (Take P to be the set of prime numbers distinct from p.)

